\subsection*{Extracting answers from pictured passages}
Can a reader choose between facts in a paragraph when the question is written by someone outside the training process? We rendered questions and passages from the Stanford Question Answering Dataset (SQuAD 1.1) as glyph images and trained the reader to return a contiguous interval of the original passage. The passage provides the language used to express the answer. Learning is required to locate it under the pictured question (Fig.~\ref{fig:optical-passage-qa}).

On 20 reserved English articles, the model returned an accepted answer under SQuAD normalization for 872 of 1,957 questions (44.6\%); mean word F1 was 59.8\%. Of these questions, 747 received the exact reference pixels including capitalization and punctuation. On the eligible subset of the official development articles, corresponding scores were 3,114 of 6,389 (48.7\%), 60.6\% F1 and 2,894 exact images. The geometry and script screen admitted 6,389 of the 10,570 development questions. In each passage with two questions having disjoint accepted answers, we took the first eligible pair in source order. The reader completed 108 of 484 reserved pairs and 351 of 1,285 development pairs. Exchanging the assignments of the two outputs gave zero and two complete pairs. On one Portugal passage, the same ink supplied the city Lisbon and the year 1822. On another paragraph, the answer ended after the United Nations when the question requested several organizations. Extractive accuracy measures one aspect of reading. The model's output support excludes answers that require composing new text.
